%preamble pfor fonts, I think Palatino is the best option for lecture notes.
\usepackage{mathpazo} % math & rm
\linespread{1.05}        % Palatino needs more leading (space between lines)
\usepackage[scaled]{helvet} % ss
\usepackage{courier} % tt
\normalfont
\usepackage[T1]{fontenc}
\usepackage{parskip}
%\usepackage{ntheorem}

%Language
\usepackage[utf8]{inputenc}
\usepackage[british]{babel}

% en



\usepackage{algorithm}
\usepackage{algpseudocode}

\usepackage{pgfplots}
\pgfplotsset{compat=1.6}
\usepackage{lmodern}

%size of the notes (FT)
\usepackage{float}
\usepackage[a4paper, portrait, margin=1.2in]{geometry}
\usepackage[font=small,labelfont=bf]{caption}
%\geometry{layoutheight=230mm,layoutwidth=160mm,layoutvoffset=30mm,layouthoffset=20mm,showcrop}
\usepackage{csquotes}


% Agregado por FVasquez
%\usepackage{tocloft}
\usepackage{setspace}

\usepackage{cancel}
\usepackage{todonotes}
\usepackage{tcolorbox}
\usepackage{hyperref}
\hypersetup{
    hidelinks = True, 
    colorlinks = false,
    linkbordercolor = {white}, 
    urlcolor = {white},
    urlbordercolor = {white}
}
\usepackage{pgfplots}
\usepackage{tikz}
\input{utils/bayesnet}
\usetikzlibrary{positioning,fit}
\input{img/week5_tools_tikz_Monge}

\usepackage{tikz}
\usetikzlibrary{calc}

% ----------------------
% Heptagram symbol
% ----------------------
\newcommand{\Heptagramold}[1][3.2cm]{%
\begin{tikzpicture}
    \def\r{#1}
    \def\offset{180/7} % = 25.714... degrees
    \pgfmathsetlengthmacro{\rin}{0.45*\r}
    \pgfmathsetlengthmacro{\rsmall}{\rin*0.5}

    % Outer circle
    %\draw[gray!60, line width=0.4pt] (0,0) circle (\r);

    % Vertices
    \foreach \i in {1,...,7}{
        \coordinate (P\i) at ({90 + (\i-1)*360/7}:\r);
        \coordinate (Q\i) at ({90 + (\i-1)*360/7}:\rin*0.69);
        \coordinate (R\i) at ({90 + \offset + (\i-1)*360/7}:\rin);
        \coordinate (S\i) at ({90 + (\i-1)*360/7}:\rsmall);
    }


    % outer heptagram
    \draw[gray!75, line width=0.9pt]
        (P1) -- (P2) -- (P3) -- (P4) -- (P5) -- (P6) -- (P7) -- cycle;

    % Inner heptagram
    \draw[gray!75, line width=0.9pt]
        (R6) -- (R1) -- (R3) -- (R5) -- (R7) -- (R2) -- (R4) -- cycle;

    % Radial lines
    \foreach \i in {1,...,7}{
        \draw[gray!75, line width=0.9pt] (P\i) -- (Q\i);
    }

    % Small heptagram
    \draw[gray!75, line width=0.9pt]
        (S1) -- (S4) -- (S7) -- (S3) -- (S6) -- (S2) -- (S5) -- cycle;

    \fill[gray!75] (0,0) circle (0.8pt);

    % Main heptagram
    \draw[line width=0.9pt]
        (P1) -- (P3) -- (P5) -- (P7) -- (P2) -- (P4) -- (P6) -- cycle;
\end{tikzpicture}%
}

\newcommand{\Heptagram}[1][3.2cm]{%
\begin{tikzpicture}
    \def\r{#1}
    \def\offset{180/7} % = 25.714... degrees
    \pgfmathsetlengthmacro{\rin}{0.9*\r}
    \pgfmathsetlengthmacro{\rsmall}{\rin*0.5}

    % Outer circle
    %\draw[gray!60, line width=0.4pt] (0,0) circle (\r);

    % Vertices
    \foreach \i in {1,...,7}{
        \coordinate (P\i) at ({90 + (\i-1)*360/7}:\r);
        \coordinate (Q\i) at ({90 + (\i-1)*360/7}:\rin);
    }


    % outer heptagram
    \draw[line width=6pt]
        (P1) -- (P3) -- (P5) -- (P7) -- (P2) -- (P4) -- (P6) -- cycle;

    % Inner heptagram
%    \draw[line width=5.7pt]
 %       (Q1) -- (Q3) -- (Q5) -- (Q7) -- (Q2) -- (Q4) -- (Q6) -- cycle;


\end{tikzpicture}%
}



\usepackage{mathtools,amssymb}
\usepackage[notocbib]{apacite}



\renewcommand{\numberline}[1]{#1.~}

%listing package para código
\usepackage{listings}
\usepackage{xcolor}
 
\definecolor{codegreen}{rgb}{0,0.6,0}
\definecolor{codegray}{rgb}{0.5,0.5,0.5}
\definecolor{codepurple}{rgb}{0.58,0,0.82}
\definecolor{backcolour}{rgb}{0.95,0.95,0.92}
 
\lstdefinestyle{mystyle}{
	xleftmargin=0.15\textwidth,
	linewidth=0.8\textwidth,
    backgroundcolor=\color{backcolour},   
    commentstyle=\color{codegreen},
    keywordstyle=\color{magenta},
    numberstyle=\tiny\color{codegray},
    stringstyle=\color{codepurple},
    basicstyle=\ttfamily\footnotesize,
    breakatwhitespace=true,         
    breaklines=true,                 
    captionpos=b,                    
    keepspaces=true,                 
    numbers=left,                    
    numbersep=5pt,                  
    showspaces=false,                
    showstringspaces=false,
    showtabs=false,                  
    tabsize=2
}
 
\lstset{style=mystyle}

\usepackage{mdframed,xcolor}

\mdfdefinestyle{TodoFrame}{%
    skipabove=\topskip,
    skipbelow=\topskip,
    topline=false,
    roundcorner=10pt,
    frametitlefont={\color{red!70!black}\normalfont\bfseries},
    frametitlerule=true,
    frametitlebackgroundcolor=red!10,
    backgroundcolor=red!10,
    linecolor=red!10,
    fontcolor=red!80!black,
}

\mdfdefinestyle{WarningFrame}{%
    skipabove=\topskip,
    skipbelow=\topskip,
    topline=false,
    roundcorner=10pt,
    frametitlefont={\color{blue!70!black}\normalfont\bfseries},
    frametitlerule=true,
    frametitlebackgroundcolor=red!10,
    backgroundcolor=red!10,
    linecolor=red!10,
    fontcolor=red!80!black,
}



\newcommand{\missing}[1]{%
  \begin{mdframed}[style=TodoFrame,nobreak=true,align=center,userdefinedwidth=0.75\textwidth]
    \textbf{To do:} #1
  \end{mdframed}
}

\newcommand{\warning}[1]{%
  \begin{mdframed}[style=WarningFrame,nobreak=true,align=center,userdefinedwidth=0.75\textwidth]
    \textbf{Warning:} #1
  \end{mdframed}
}




\usepackage{amsmath}
%theorems
\usepackage{amsthm}
\usepackage{amssymb}

\renewcommand\qedsymbol{$\blacksquare$}

 


% Teoremas

% Agregada por FVasquez
\newtheoremstyle{mytheoremstyle} % name
    {8pt}                % Space above
    {8pt}                % Space below
    {}                   % Body font
    {20pt}                   % Indent amount
    {\bfseries}          % Theorem head font
    {:}                  % Punctuation after theorem head
    {\newline}               % Space after theorem head
    {}                   % Theorem head spec
% Agregada por FVasquez

\newtheoremstyle{mytheoremstyle}
{.5\topsep}
{.5\topsep}
{\addtolength{\leftskip}{2.5em}}
{-2.5em}
{\bfseries}
{.}
{\newline}
{}

\theoremstyle{mytheoremstyle}

\newtheorem{definition}{Definition}[chapter]
\newtheorem{example}{Example}[chapter]
\newtheorem{exercise}{Exercise}[chapter]
\newtheorem{solution}{Solution}[chapter]
\newtheorem{theorem}{Theorem}[chapter]
\newtheorem{proposition}{Proposition}[chapter]
\newtheorem{lemma}{Lemma}[chapter]
\newtheorem{remark}{Remark}[chapter]
\newtheorem{corollary}{Corollary}


%macros
\newcommand*\clase[1]{\vspace{2em}{\color{red}\par\noindent\raisebox{.8ex}{\makebox[\linewidth]{\hrulefill\hspace{1ex}\raisebox{-.8ex}{#1}\hspace{1ex}\hrulefill}}\vspace{0em}}}

\usepackage{color}
\providecommand{\red}[1]{\textcolor{red}{{\bf #1}}}


% Agregado por FVasquez
\usepackage{titlesec}
\titleformat*{\section}{\LARGE\bfseries}
\titleformat*{\subsection}{\Large\bfseries}
\titleformat*{\subsubsection}{\large\bfseries}
\titleformat*{\paragraph}{\large\bfseries}
\titleformat*{\subparagraph}{\large\bfseries}
\titleformat{\chapter}[display]
  {\normalfont\huge\bfseries\filcenter}{Week \ \thechapter}{20pt}{\Huge}
\titlespacing*{\chapter}
  {0pt}{30pt}{20pt}

\def\familiaparametrica{\mathcal{P}  = \{P_\theta\ \tq\ \theta\in\Theta\}}
\def\familiaparametricaO{\mathcal{P}  = \{P_\theta\ \tq\ \theta\in\Omega\}}
\def\densidadparametrica{\{p_\theta\ \tq\ \theta\in\Theta\}}
\newcommand{\E}[1]{\mathbb{E} \left(#1\right)}
\newcommand{\V}[1]{\mathbb{V} \left(#1\right)}
\newcommand{\var}[1]{\text{var} \left(#1\right)}
\newcommand{\mean}[1]{\text{mean} \left(#1\right)}
\newcommand{\Prob}[1]{\mathbb{P} \left(#1\right)}
\newcommand{\PP}{\mathbb{P}}
\newcommand{\Et}[1]{\mathbb{E}_\theta \left(#1\right)}
\newcommand{\Vt}[1]{\mathbb{V}_\theta \left(#1\right)}
\newcommand{\Probt}[1]{\mathbb{P}_\theta \left(#1\right)}
\newcommand{\Probtu}[1]{\mathbb{P}_{\theta_1} \left(#1\right)}
\newcommand{\Probtz}[1]{\mathbb{P}_{\theta_0} \left(#1\right)}
\newcommand{\cat}[1]{\text{Cat}\left(#1\right)}
\newcommand{\dir}[1]{\text{Dir}\left(#1\right)}
\newcommand{\ber}[1]{\text{Ber}\left(#1\right)}
\newcommand{\bet}[1]{\text{Beta}\left(#1\right)}
\newcommand{\dete}[1]{\text{det} #1 }
\newcommand{\bin}[1]{\text{Bin}\left(#1\right)}
\newcommand{\mul}[1]{\text{Mult}\left(#1\right)}
\newcommand{\uni}[1]{\text{Uniform}\left(#1\right)}
\newcommand{\poi}[1]{\text{Poisson}\left(#1\right)}
\newcommand{\expo}[1]{\exp\left(#1\right)}
\newcommand{\loga}[1]{\log\left(#1\right)}
\newcommand{\KL}[2]{\text{KL}\left(#1\middle\|#2\right)}
\newcommand{\DL}[2]{L_1\left(#1\middle\|#2\right)}
\newcommand{\Dchi}[2]{\chi^2\left(#1\middle\|#2\right)}
\newcommand{\one}[1]{\mathbb{1}_{#1}}
\newcommand{\defeq}{\mathrel{\stackrel{\mathrm{def}}{=}}}
\newcommand{\GP}[2]{\mathcal{GP}\left(#1 , #2\right)}
\newcommand{\tr}{\operatorname{tr}}
\newcommand{\indep}{\perp \!\!\! \perp}


\def\t{{\mathbf t}}
\def\cA{{\mathcal A}}
\def\cB{{\mathcal B}}
\def\cC{{\mathcal C}}
\def\cD{{\mathcal D}}
\def\cF{{\mathcal F}}
\def\cL{{\mathcal L}}
\def\cN{{\mathcal N}}
\def\cO{{\mathcal O}}
\def\cP{{\mathcal P}}
\def\cQ{{\mathcal Q}}
\def\cT{{\mathcal T}}
\def\cX{{\mathcal X}}
\def\cY{{\mathcal Y}}
\def\cZ{{\mathcal Z}}

\def\d{{\text{d}}}
\def\diag{\mathop{\rm diag}\nolimits}%
\def\dt{{\text{d}t}}
\def\dtheta{{\text{d}\theta}}
\def\dx{{\text{d}x}}
\def\dy{{\text{d}y}}
\def\dz{{\text{d}z}}

\def\ee{{\text{ee}}}
\def\elbo{{\text{ELBO}}}

\def\f{{ \mathbf{f}}}
\def\gh{{ \hat{g}}}
\def\ghx{{ \hat{g}(x)}}
\def\ghX{{ \hat{g}(X)}}

\def\ind{{\mathbb 1}}
\def\I{{\mathbf I}}
\def\m{{ \mathbf{m}}}

\def\N{{\mathbb N}}
\def\one{{ \mathbb{1}}}
\def\P{{\mathbb P}}

\def\R{{\mathbb R}}

\def\thetaMV{{\theta_\text{MV}}}
\def\thetahat{{\hat\theta}}
\def\tq{{\text{ t.q. }}}

\def\u{{ \mathbf{u}}}
\def\v{{ \mathbf{v}}}
\def\w{{ \mathbf{w}}}

\def\x{{ \mathbf{x}}}
\def\X{{ \mathbf{X}}}
\def\xb{{ \bar{x}}}

\def\y{{ \mathbf{y}}}
\def\z{{ \mathbf{z}}}

\DeclareMathOperator*{\argmax}{arg\,max}
\DeclareMathOperator*{\argmin}{arg\,min}